\documentclass[12pt]{article}
\usepackage[utf8]{inputenc}
\usepackage[T1]{fontenc}
\usepackage{hyperref}
\title{Turiw\'ara Language Resources for K-12 Curriculum Development}
\author{Compiled by Claude Sonnet 5 (claude-sonnet-5) for Andr\'as Kornai}
\date{}
\begin{document}
\maketitle
\section{Introduction}
Turiw\'ara (ISO 639-3 twt) is a Tupi-Guaran\'i language of Par\'a, Brazil, probably no longer spoken. The ASJP database gives it as extinct since 1990, while Joshua Project lists 60 speakers, probably from older data. A reported 1995 population was about 30.
\section{Primary Educational Resources}
None found.
\section{Linguistic and Reference Materials}
\begin{itemize}
\item Wikipedia on the language: \url{https://en.wikipedia.org/wiki/Turiw%C3%A1ra_language}
\item Wikipedia on the people: \url{https://en.wikipedia.org/wiki/Turiw%C3%A1ra}
\item ASJP word list: \url{https://asjp.clld.org/languages/TURIWARA}
\item Glottolog entry (parent languoid): \url{https://glottolog.org/resource/languoid/id/para1312}
\end{itemize}
\section{Religious and Community Resources}
Joshua Project people-group page: \url{https://joshuaproject.net/people_groups/21568/BR}
\section{Digital and Audio Resources}
None found.
\section{Curriculum Development Context}
Documentation is a short word list at most. A heritage curriculum would have to be about the people and history, with at most a few words of the language, and would need to be developed with the community.
\end{document}